\section{反標準環からの導来圏の復元：正確な判定と障害}
\label{sec:restored-anticanonical}
第\ref{sec:anticanonical-common-ring}節は、反標準 orbit 環を
二点 blow-up 環と同定する。しかし環の同型とカテゴリーの同値は別の命題である。
その間に必要な条件を、既知の復元定理を用いて明示する。

\begin{theorem}[反標準 section functor の十分条件]\label{thm:restored-az}
$\mathcal H$ を $k$-線形、局所 noetherian な abelian category とし、
$O\in\mathcal H$ を noetherian object、$F:\mathcal H\to\mathcal H$ を
完全な自己同値とする。
任意の noetherian objects の間の Hom が有限次元であり、
$(O,F)$ が Artin--Zhang の意味で ample、すなわち
\begin{enumerate}
\item 任意の noetherian な $M\in\mathcal H$ に対し、有限個の $n_i\ge0$ と全射
$\bigoplus_iF^{-n_i}O\twoheadrightarrow M$ が存在する。
\item 任意の noetherian objects の全射 $M\twoheadrightarrow N$ について、
$n\gg0$ のとき
$\operatorname{Hom}(F^{-n}O,M)\to\operatorname{Hom}(F^{-n}O,N)$ が全射となる。
\end{enumerate}
と仮定する。
$R=\bigoplus_{n\ge0}\operatorname{Hom}_{\mathcal H}(O,F^nO)$ に
$u\cdot v=F^{\deg v}(u)\circ v$ で積を入れると、
$R$ は右 noetherian であり、自然な section functor は
$\mathcal H\simeq\operatorname{qgr}R$ を誘導する。
さらに、ある smooth proper dg 圏 $\mathcal C$ が
$D^b(\mathcal H)$ の指定 dg enhancement に同値ならば、
$\mathcal C\simeq D^b(\operatorname{qgr}R)$ である。
\end{theorem}
\begin{proof}
最初の結論は Artin--Zhang の ample pair による復元定理である
\cite{KeelerNCRings,ArtinZhang1994}。
後半は同じ abelian equivalence を有界導来圏へ適用する。
この定理は $F$ が \emph{指定した heart 上で t-exact} であることを前提とする。
三角圏上の自己同値であるというだけでは仮定を満たさない。
\end{proof}

\begin{corollary}[本研究に必要な検証項目]\label{cor:restored-az-criteria}
第\ref{sec:anticanonical-common-ring}節の階数23の共通圏 $\mathcal C_B$
について、反標準関手 $F=S_{\mathcal C_B}^{-1}[2]$ と
構造対象 $O$ が、ある noetherian な bounded heart $\mathcal H_B$ に対して
$F$ の t-exactness と定理\ref{thm:restored-az}の二つの ample 条件を満たし、
$D^b(\mathcal H_B)\simeq\mathcal C_B$ が成り立てば、
\[
\mathcal C_B\simeq D^b(\operatorname{qgr}\mathscr A),
\qquad
\mathscr A=\bigoplus_{n\ge0}\operatorname{Hom}(O,F^nO)
\]
となる。特に固定された無限位数ファイバーにおける
$\mathscr A\simeq\mathcal S^{(3)}(\tau^{-1}p+\tau^{-1}q)$ の
\emph{環としての}同定からこの結論を得るには、上記の heart と ample 性の検証が必要である。
\end{corollary}

\begin{proposition}[五つの section module の比較の十分条件]
$R$ を右 noetherian な次数付き代数とし、
$\mathcal C_B=\operatorname{Perf}(B)$ に五個の full exceptional objects
$E_1,\ldots,E_5$ を選ぶ。
$\operatorname{qgr}R$ 上の五個の objects $M_i$ が、
\[
\operatorname{RHom}_{\operatorname{qgr}R}(M_i,M_j)
 \simeq \operatorname{RHom}_{\mathcal C_B}(E_i,E_j)
\]
を合成と単位まで保つ dg 代数として実現し、
$\bigoplus_iM_i$ が $D^b(\operatorname{qgr}R)$ の classical generator
であるとする。このとき
$D^b(\operatorname{qgr}R)\simeq\mathcal C_B$ が成立する。
\end{proposition}
\begin{proof}
五個の対象の full dg endomorphism algebra を比較し、
導来 Morita 理論を適用する。
単に Euler 行列や $K_0$ 上の Serre 作用が一致するだけでは、
full dg endomorphism algebra の一致は従わない。
\end{proof}

\begin{remark}[再開時の有限入力]
もし定理\ref{thm:restored-az}の heart が得られれば、
各 $E_i$ に対して
\[
M_i=\bigoplus_{n\ge0}\operatorname{Hom}_{\mathcal H_B}(O,F^nE_i)
\]
を考えられる。$u\in\operatorname{Hom}(O,F^nE_i)$ と
$r\in R_m$ の右積は $u\cdot r=F^m(u)\circ r$ である。
十分大きな共通の次数シフトを選び、五つの加群の有限表示、
全高次 $\operatorname{Ext}$、生成性、各積 $M_i\otimes M_j\to M_k$
に相当する合成の比較を調べればよい。
第\ref{sec:anticanonical-common-ring}節に既にある
$\mathcal S_3$ の二点 evaluation の核と42本の二次関係式は、
この有限計算の入力となる。
ただし\emph{すべての次数についての生成性}と
$F$ の t-exactness は現時点で検証されていない。
\end{remark}

\begin{remark}[反標準環だけでは圏を復元できない例]\label{ex:restored-weak-dp}
$Y$ を $\mathbb P^2$ の一点 $p$ を blow-up し、
さらに最初の例外曲線上の点 $q$ を blow-up した滑らかな曲面とする。
$H,E_1,E_2$ を pullback basis とすると
\[
 -K_Y=3H-E_1-E_2,\qquad C=E_1-E_2,
 \quad C^2=-2,\quad (-K_Y)\cdot C=0.
\]
ここで $C$ は最初の例外曲線の狭義変換である。
$-K_Y$ は big かつ nef だが ample ではなく、反標準射は $C$ を収縮する。
反標準環 $R(Y,-K_Y)$ の $\operatorname{Proj}$ は $A_1$ 型特異点を持つ
反標準模型になる。従って
$\operatorname{coh}Y$ を $\operatorname{qgr}R(Y,-K_Y)$ と
反標準 section functor で同一視することはできない。
実際、右辺は特異な反標準模型の coherent sheaves であり、
$Y$ の滑らかな幾何をそのまま回復しない。
これは固定ファイバーの環の同型だけから圏の同値を
推論できないことを示す可換な障害例である。
\end{remark}
